\chapter{LLM agents and tool calling}
\label{ch:agents}

An \term{LLM agent} is a loop: the model observes state, decides on an action (usually a
tool call), the runtime executes it, and the observation feeds back in. Everything the
decision layer does hangs off this loop, so it is worth being precise about it.

\section{The agent loop}

\begin{figure}[h]\centering
\begin{tikzpicture}[node distance=7mm and 14mm]
  \node[frontend] (obs) {observe state};
  \node[model, right=of obs] (decide) {LLM decides\\next action};
  \node[backend, right=of decide] (exec) {runtime executes\\tool};
  \node[infra, below=8mm of decide] (env) {environment /\\tools};
  \draw[flow] (obs) -- (decide);
  \draw[flow] (decide) -- node[flowlabel]{tool call} (exec);
  \draw[flow] (exec) -- (env);
  \draw[flow] (env.west) to[bend left=20] node[flowlabel]{result} (obs.south);
\end{tikzpicture}
\caption{The agent loop. The decision layer intercepts the ``LLM decides'' step for
bounded, closed-set decisions.\datalabel{Illustrative}}
\label{fig:agentloop}
\end{figure}

\section{Tool calling guarantees shape, not confidence}

Modern APIs let the model emit a structured \term{tool call} (a function name plus JSON
arguments). Structured outputs and strict tool use guarantee the \emph{shape} is valid, but
hosted APIs generally do not return a probability distribution over \emph{which} tool was
chosen. That missing distribution is exactly the primitive the decision layer supplies.

\begin{keyidea}
Frameworks already expose a ``pre-action'' seam: a callable that returns an enum
(allow/deny/ask) or a node name. Escalation there is a boolean, never a probability. The
decision layer plugs into that seam and makes the signal a calibrated probability instead.
\end{keyidea}

\section{Where decisions live at scale}

When a catalog has dozens or hundreds of tools, accuracy degrades: vendors themselves note
tool selection gets worse past roughly 30--50 tools, which is why they ship tool-search and
allowed-tools features. Narrowing the catalog \emph{before} the LLM chooses is a real job,
and a cheap classifier is a natural fit for it. That is decision point DP2, the MVP's first
target (Chapter~\ref{ch:decisionpoints}).

\begin{pitfall}
A tempting shortcut is to let the model both narrow and choose in one call. That reintroduces
the cost you were trying to avoid and hides the decision inside free text, where you cannot
attach a confidence or test it. Keep the bounded decision separate and typed.
\end{pitfall}

\begin{checkbox}
Name two things a hosted tool-calling API guarantees and one thing it does not. Which of the
three does the decision layer add? (Appendix~\ref{app:answers}.)
\end{checkbox}
